\documentclass[10pt,a4paper]{amsart}
\usepackage[margin=19mm]{geometry}
\usepackage{amsmath,amssymb,mathtools}
\usepackage[scr=rsfs,cal=euler]{mathalfa}
\usepackage{xfrac}
\usepackage[unicode,hidelinks,bookmarks=false]{hyperref}
\newcommand{\ie}{i.e.\ }
\newcommand{\cf}{cf.\ }
\newcommand{\resp}{respectively\ }
\newcommand{\Subsection}{\subsection}
\providecommand{\nobreakdash}{\nobreak\hbox{-}}
\setlength{\emergencystretch}{2em}
\multlinegap0pt
\usepackage{bm}
\usepackage{bbm}
\usepackage{dsfont}
\usepackage{xspace}
\usepackage{cancel}
\usepackage{mathtools}
\usepackage{amsthm}
\makeatletter
\@ifundefined{Theorem}{\newtheorem{Theorem}{Theorem}}{}
\makeatother
\title[A Tensor Category Construction of the $W\_{p,q}$ Triplet Vert]{A Tensor Category Construction of the $W\_{p,q}$ Triplet Vertex Operator Algebra and Applications}
\author{Robert McRae, Valerii Sopin}
\date{Source: arXiv:2508.18895v2 (CC BY 4.0)}
\begin{document}
\maketitle

\noindent\emph{Source and licence.} A Tensor Category Construction of the $W_{p,q}$ Triplet Vertex Operator Algebra and Applications. Version arXiv:2508.18895v2 (CC BY 4.0), distributed under the Creative Commons Attribution 4.0 International licence, as recorded in the arXiv licence metadata for this exact version and shown on the abstract page https://arxiv.org/abs/2508.18895v2. Full rights notice: licenses/arxiv-2508.18895-CC-BY-4.0.txt.\par\smallskip

\noindent\textit{Editorial scope.} Counted unit: the Theorem whose source label is \texttt{thm:C'-alg-to-C-alg} in the article (source0.tex lines 773--780, source label \texttt{thm:C'-alg-to-C-alg}), delivered together with the article's own complete original proof of it (source0.tex lines 782--836, one \texttt{proof} environment of 55 lines). No mathematics is written, repaired or re-derived: statement and proof are the article's own text line for line. Required context retained uncounted: eqn:l'-r'-gen-def (lines 733--735); eqn:A0-iso (lines 740--742); eqn:An-iso (lines 744--746); eqn:A'-A-structure-reln (lines 759--761); eqn:g'-g-reln (lines 768--770). Every definition, display and label of those lines is retained unchanged. Omitted from this package and not redistributed: 1--732, 736--739, 743--743, 747--758, 762--767, 771--772, 781--781, 837--1780. Nothing inside a retained excerpt is omitted or altered: every retained body is delivered exactly as the article's lines are written, so no omission receipt of any kind accompanies this package. Citations: the retained excerpts carry no citation command at all, so no bibliography entry of the article is transcribed and no reference is redistributed. Reference adaptation: none was needed and none was made; every reference inside the delivered unit resolves inside it, so no unresolved reference or citation is produced. Printed slips kept literal: no printed word, symbol, hypothesis or number of the counted statement or of its proof was altered. Layout and class: the article's own class is \texttt{sigma} with packages including ifpdf, mathtools, xy, tikz, tikz-cd; the wrapper is the lane's pinned \texttt{amsart} editorial wrapper at 10pt on A4, which restates the theorem-like environments the retained text needs and adds the article's own macro definitions that the retained text needs (none), with extra wrapper packages (bm, bbm, dsfont, xspace, cancel, mathtools). The delivered numbering is the unit's own. Assets: no retained span contains a figure environment, a graphics-inclusion command, a PDF-page inclusion, a table or a TikZ picture; no image file is copied into this package, the package declares no asset dependency and no image file is redistributed with it. Licence: version arXiv:2508.18895v2 of this article is distributed under the Creative Commons Attribution 4.0 International licence, as recorded in the arXiv licence metadata for this exact version and shown on the abstract page https://arxiv.org/abs/2508.18895v2. A full rights notice is delivered at licenses/arxiv-2508.18895-CC-BY-4.0.txt. Characters: every retained excerpt is pure ASCII, so the delivered engine, XeLaTeX with its default Latin Modern text font, reports no missing character.
\begin{gather}\label{eqn:l'-r'-gen-def}
l_X' = l_X\circ(\varphi\boxtimes\mathrm{Id}_X),\qquad r_X'=r_X\circ(\mathrm{Id}_X\boxtimes\varphi)
\end{gather}

\begin{gather}\label{eqn:A0-iso}
\Phi\circ -\colon\ \operatorname{Hom}(\mathbf{1}',A')\longrightarrow\operatorname{Hom}(\mathbf{1}',A),\qquad -\circ\varphi\colon\ \operatorname{Hom}(\mathbf{1},A)\longrightarrow\operatorname{Hom}(\mathbf{1}',A)
\end{gather}

\begin{gather}
\Phi\circ -\colon\ \operatorname{Hom}\bigl(\bigl(A'\bigr)^{\boxtimes n},A'\bigr)\longrightarrow\operatorname{Hom}\bigl(\bigl(A'\bigr)^{\boxtimes n},A\bigr),\nonumber\\ -\circ\Phi^{\boxtimes n}\colon\ \operatorname{Hom}\bigl(A^{\boxtimes n},A\bigr)\longrightarrow\operatorname{Hom}\bigl(\bigl(A'\bigr)^{\boxtimes n},A\bigr)\label{eqn:An-iso}
\end{gather}

\begin{gather}\label{eqn:A'-A-structure-reln}
\Phi\circ\iota_{A'} =\iota_A\circ\varphi,\qquad \Phi\circ\mu_{A'} =\mu_A\circ(\Phi\boxtimes\Phi).
\end{gather}

\begin{gather}\label{eqn:g'-g-reln}
\Phi\circ g' = g\circ\Phi.
\end{gather}

\begin{Theorem}\label{thm:C'-alg-to-C-alg}
In the setting of the previous paragraph,
\begin{enumerate}\itemsep=0pt
\item[$(1)$] $(A',\mu_{A'},\iota_{A'})$ is a commutative algebra in $\mathcal{C}'$ if and only if $(A,\mu_A,\iota_A)$ is a commutative algebra in $\mathcal{C}$, where $(\iota_{A'},\mu_{A'})$ and $(\iota_A,\mu_A)$ are related by \eqref{eqn:A'-A-structure-reln}.

\item[$(2)$] A $\mathcal{C}'$-morphism $g'\colon A'\rightarrow A'$ is an algebra isomorphism between two commutative algebra structures \smash{$\bigl(A',\mu_{A'}^{(1)},\iota_{A'}^{(1)}\bigr)$} and \smash{$\bigl(A',\mu_{A'}^{(2)},\iota_{A'}^{(2)}\bigr)$} if and only if $g\colon A\rightarrow A$ defined by \eqref{eqn:g'-g-reln} is an~algebra isomorphism between \smash{$\bigl(A,\mu_{A}^{(1)},\iota_{A}^{(1)}\bigr)$} and \smash{$\bigl(A,\mu_{A}^{(2)},\iota_{A}^{(2)}\bigr)$}, where \smash{$\bigl(\mu_A^{(i)},\iota_A^{(i)}\bigr)$} for $i=1,2$ are defined by \eqref{eqn:A'-A-structure-reln}.
\end{enumerate}
\end{Theorem}

\begin{proof}
(1) For the left unit property of a commutative algebra, we calculate
\begin{align*}
\mu_A\circ(\iota_A\boxtimes\mathrm{Id}_A)\circ l_A^{-1}\circ\Phi & = \mu_A\circ(\iota_A\boxtimes\mathrm{Id}_A)\circ(\mathrm{Id}_{\mathbf{1}}\boxtimes\Phi)\circ l_{A'}^{-1}\nonumber\\
& =\mu_A\circ(\mathrm{Id}_{A}\boxtimes\Phi)\circ(\iota_A\boxtimes\mathrm{Id}_{A'})\circ(\varphi\boxtimes\mathrm{Id}_{A'})\circ(l_{A'}')^{-1}\nonumber\\
& =\mu_A\circ(\Phi\boxtimes\Phi)\circ(\iota_{A'}\boxtimes\mathrm{Id}_{A'})\circ(l_{A'}')^{-1} \nonumber\\
&=\Phi\circ\mu_{A'}\circ(\iota_{A'}\boxtimes\mathrm{Id}_{A'})\circ(l_{A'}')^{-1}
\end{align*}
using the naturality of the left unit isomorphisms in $\mathcal{C}$, \eqref{eqn:l'-r'-gen-def}, and \eqref{eqn:A'-A-structure-reln}.
Since the two maps in~\eqref{eqn:An-iso} are isomorphisms in the $n=1$ case, it follows that
\begin{gather*}
\mu_A\circ(\iota_A\boxtimes\mathrm{Id}_A)\circ l_A^{-1} =\mathrm{Id}_A\longleftrightarrow \mu_{A'}\circ(\iota_{A'}\boxtimes\mathrm{Id}_{A'})\circ(l_{A'}')^{-1}=\mathrm{Id}_{A'}.
\end{gather*}
Thus the left unit property holds for $(A',\mu_{A'},\iota_{A'})$ if and only if it holds for $(A,\mu_A,\iota_A)$.

It follows similarly that the right unit property for $(A',\mu_{A'},\iota_{A'})$ is equivalent to the right unit property for $(A,\mu_A,\iota_A)$. Alternatively, this follows from the left unit property and the equivalence of the commutativity of $\mu_{A'}$ and $\mu_A$, which we prove next. Since
\begin{align*}
\mu_A\circ\mathcal{R}_{A,A}\circ(\Phi\boxtimes\Phi) =\Phi\circ\mu_{A'}\circ\mathcal{R}_{A',A'},\qquad\mu_A\circ(\Phi\boxtimes\Phi)=\Phi\circ\mu_{A'}
\end{align*}
by \eqref{eqn:A'-A-structure-reln} and the naturality of the braiding isomorphisms, the assumption that the two maps in~\eqref{eqn:An-iso} are isomorphisms in the $n=2$ case implies that $\mu_A\circ\mathcal{R}_{A,A}=\mu_A$ if and only if $\mu_{A'}\circ\mathcal{R}_{A',A'}=\mu_{A'}$. Thus $\mu_A$ is commutative if and only if $\mu_{A'}$ is. Similarly, $\mu_A$ is associative if and only if $\mu_{A'}$ is because
\begin{align*}
\mu_A\circ(\mathrm{Id}_A\boxtimes\mu_A)\circ(\Phi\boxtimes(\Phi\boxtimes\Phi)) =\Phi\circ\mu_{A'}\circ(\mathrm{Id}_{A'}\boxtimes\mu_{A'})
\end{align*}
and
\begin{gather*}
\mu_A\circ(\mu_A\boxtimes\mathrm{Id}_A)\circ\mathcal{A}_{A,A,A}\circ(\Phi\boxtimes(\Phi\boxtimes\Phi)) = \Phi\circ\mu_{A'}\circ(\mu_{A'}\boxtimes\mathrm{Id}_{A'})\circ\mathcal{A}_{A',A',A'}
\end{gather*}
by \eqref{eqn:A'-A-structure-reln} and naturality of the associativity isomorphisms, and because the two maps in \eqref{eqn:An-iso} are isomorphisms in the $n=3$ case.
This proves the first statement of the theorem.

(2) Now suppose $g'\colon A'\rightarrow A'$ and $g\colon A\rightarrow A$ are two morphisms related by~\eqref{eqn:g'-g-reln}. If $g$ is a~$\mathcal{C}$-isomorphism with inverse $g^{-1}$, then $g'$ is a $\mathcal{C}'$-isomorphism with inverse $\bigl(g^{-1}\bigr)'$ characterized by \smash{$\Phi\circ\bigl(g^{-1}\bigr)'=g^{-1}\circ\Phi$} as in \eqref{eqn:g'-g-reln}. Indeed,
\begin{gather*}
\Phi\circ g'\circ\bigl(g^{-1}\bigr)' =g\circ g^{-1}\circ\Phi =\Phi
\end{gather*}
by \eqref{eqn:g'-g-reln}, and therefore \smash{$g'\circ\bigl(g^{-1}\bigr)'=\mathrm{Id}_{A'}$} since the first map of \eqref{eqn:An-iso} is an isomorphism in the $n=1$ case. Similarly, \smash{$\bigl(g^{-1}\bigr)'\circ g' =\mathrm{Id}_{A'}$}, and similarly $g$ is a~$\mathcal{C}$-isomorphism if $g'$ is a~$\mathcal{C}'$-isomorphism.

Now we consider how $g$ and $g'$ relate to different algebra structures on $A$ and $A'$. We have
\begin{gather*}
g\circ\iota^{(1)}_A\circ\varphi =\Phi\circ g'\circ\iota_{A'}^{(1)},\qquad \iota_{A}^{(2)}\circ\varphi =\Phi\circ\iota_{A'}^{(2)}
\end{gather*}
and
\begin{gather*}
g\circ\mu_A^{(1)}\circ(\Phi\boxtimes\Phi) =\Phi\circ g'\circ\mu_{A'}^{(1)},\qquad \mu_{A}^{(2)}\circ(g\boxtimes g)\circ(\Phi\boxtimes\Phi) =\Phi\circ\mu_{A'}^{(2)}\circ(g'\boxtimes g')
\end{gather*}
by \eqref{eqn:A'-A-structure-reln} and \eqref{eqn:g'-g-reln}. Thus
\begin{gather*}
g\circ\iota_A^{(1)}=\iota_A^{(2)}\longleftrightarrow g'\circ\iota_{A'}^{(1)} =\iota_{A'}^{(2)}
\end{gather*}
since the maps in \eqref{eqn:A0-iso} are an isomorphism, and
\begin{gather*}
g\circ\mu_A^{(1)}=\mu_A^{(2)}\circ(g\boxtimes g) \longleftrightarrow g'\circ\mu_{A'}^{(1)}=\mu_{A'}^{(2)}\circ(g'\boxtimes g')
\end{gather*}
since the maps in \eqref{eqn:An-iso} are isomorphisms in the $n=2$ case.
This proves that $g\colon \smash{\bigl(A,\mu\raisebox{1pt}{${}_A^{(1)}$},\iota\raisebox{1pt}{${}_A^{(1)}$}\bigr)}\!\rightarrow \smash{\bigl(A,\mu\raisebox{-1pt}{${}_A^{(2)}$}, \iota\raisebox{-1pt}{${}_A^{(2)}$}\bigr)}$ is an isomorphism of $\mathcal{C}$-algebras if and only if \smash{$g'\colon \bigl(A',\mu\raisebox{-1pt}{${}_{A'}^{(1)}$},\iota\raisebox{-1pt}{${}_{A'}^{(1)}$}\bigr) \!\rightarrow\!\bigl(A',\mu\raisebox{-1pt}{${}_{A'}^{(2)}$},\iota\raisebox{-1pt}{${}_{A'}^{(2)}$}\bigr)$} is an~isomorphism of $\mathcal{C}'$-algebras.
\end{proof}

\end{document}
